% !TEX root = ../../Groups and Monoids in Context.tex
\section{Tuples and maps}
Now we have a notion of set, we should talk about what sorts of objects will frequently appear besides numbers.
\begin{definition}[products and $n$-tuples]
    If we have a set $A$ and another set $B$, we can form the (cartesian) product of $A$ with $B$:
    \begin{equation*}
        A \times B := \{(a,b) \mid a \in A, b \in B \}.
    \end{equation*}
    Here, $(a,b)$ is an ordered pair, so $(a,b) = (c,d)$ if and only if $a = c$ and $b = d$.

    \medskip\noindent
    Similarly, if we have $n\in \N$ and $n$ sets $A_1,\dots,A_n$, we can form their cartesian product:
    \begin{equation*}
        \prod_{i=1}^n A_i := \{ (a_1,\dots,a_n) \mid \text{$a_i \in A_i$ for each $i \in \llb 1,n \rrb$}\}.
    \end{equation*}
    where each $(a_1,\dots,a_n)$ is an $n$-tuple which is again ordered, so $(a_1,\dots,a_n) = (b_1,\dots,b_n)$ if and only if $a_i = b_i$ for every $i \in \llb 1,n \rrb$.

    \medskip\noindent
    In the case where all $A = A_i$ for every $i \in \llb 1, n \rrb$, we may also denote this set as $A^n$.
\end{definition}
\begin{remark}
    Ordered pairs are just another name for 2-tuples. 3-tuples are often called triples. It is acceptable (but rare) to continue using this naming scheme.
\end{remark}
\begin{remark}
    Notice that for $a,b \in \R$ with $a < b$, the expression $(a,b)$ is ambiguous: it could refer to an open interval, and it could refer to an ordered pair.
    These are unfortunately, both fairly standard notations, and thus, this ambiguity appears a lot.
    It does not generally cause much confusion though, since context almost always makes it clear which way it should be interpreted\footnote{In programming, objects often have a specific type. The same concept can be used to formalise mathematics. This approach largely solves the problem of ambiguities like this, since $(a,b)$ will be specified as either of type Interval or type $\R^2$.}.
\end{remark}

\noindent
Tuples are normally used to define functions, but the definition which uses them is not a helpful definition, it is simply one made out of necessity when people were formalising mathematics with set theory. As such, I present an alternative definition:
\begin{definition}[Maps]
    A map (or function) $f$ is an object with a domain set $A$ and a codomain $B$ which, to each element of $A$ assigns a unique element of $B$. This is denoted $f \colon A \to B$, read ``$f$ is a map from $A$ to $B$''.

    \medskip\noindent
    The unique element of $B$ assigned to a specific $a \in A$ is written $f(a)$, read ``$f$ of $a$'', so $f(a) \in B$ for any $a \in A$.

    \medskip\noindent
    We can also write $f\colon a \mapsto b$ for $a\in A, b\in B$ to say that $f(a) = b$.

    Two functions $f,g \colon A \to B$ are considered equal if $f(a) = g(a)$ for every $a \in A$.
\end{definition}
\begin{remark}
    The unique assignment a function performs does not have to be described by an algorithm to produce it in general.
\end{remark}
\begin{example}
    Following on from the remark, the following is a valid function $\llb 1, 5 \rrb \to \{\text{\emoji{grapes}},\text{\emoji{cat}},\text{\emoji{house}}\}$:
    \begin{align*}
                 1 & \longmapsto \text{\emoji{grapes}} \\
                 2 & \longmapsto \text{\emoji{house}} \\
        f \colon 3 & \longmapsto \text{\emoji{house}} \\
                 4 & \longmapsto \text{\emoji{grapes}} \\
                 5 & \longmapsto \text{\emoji{house}}
    \end{align*}
\end{example}
\begin{example}[Identity]
    For any set $S$, there is a function $\Id_S \colon S \to S$ which does nothing:
    \begin{equation*}
        \Id_S \colon s \mapsto s
    \end{equation*}
    so, $\Id_S(s) = s$.
\end{example}

\begin{definition}[Composition of functions]
    A pair of functions $f \colon A \to B$, $g \colon B \to C$ can be composed to obtain a function $g \circ f \colon A \to C$ given by:
    \begin{equation*}
        g \circ f \colon a \overset{f}{\longmapsto} f(a) \overset{g}{\longmapsto} g(f(a))
    \end{equation*}
\end{definition}
\begin{remark}
    It's important to notice that $g \circ f$ does not mean ``apply $g$, then apply $f$'' but instead the reverse.
    This can confuse people, but this convention is so that the order of functions matches how you would write application without the composition operation:
    \begin{equation*}
        (g \circ f)(a) = g(f(a))
    \end{equation*}
    This is largely due to the convention of how function application is written\footnote{Alternatives conventions include: $f x$, $(x)f$, $x f$. ie. omitting the brackets when unnecessary, reversing the order, and the combination of both. Arguably the last of these is the most natural to use in the contexts this book will cover, but I have chosen the most common notation for reader familiarity.}.
\end{remark}

\begin{definition}[Injections, surjection, and bijections]
    A map $f \colon A \to B$ is called injective if, whenever $f(a_1) = f(a_2)$, $a_1 = a_2$. That is, no two elements of $A$ are mapped to the same element of $B$. Such a map is called an injection.

    \medskip\noindent
    A map $f \colon A \to B$ is called surjective if, for every $b \in B$, there is some $a \in A$ with $f(a) = b$. That is, $f$ maps onto every element of $B$. Such a map is called a surjection.

    \medskip\noindent
    A map $f \colon A \to B$ is called bijective if it is both injective and surjective.
\end{definition}

\begin{definition}
    A map $f \colon A \to B$ is called invertible if there is a map $g \colon B \to A$ such that $g \circ f = \Id_A$ and $f \circ g = \Id_B$. The map $g$ is called the inverse of $f$.
\end{definition}
\begin{remark}
    Formally, this definition does not prove that a function only has 1 inverse if it has one, but Exercise \ref{unique_inverse} shows that inverses are unique if they exist.
\end{remark}

\begin{example}[Identity maps are invertible and bijective]
    To show $\Id_S \colon S \to S$ is invertible, we must find an inverse function $f \colon S \to S$ which reverses the operation of $\Id_S$. But we said that the identity maps do nothing to their elements, so surely if we do nothing to them twice we'll still end up doing nothing? Indeed,
    \begin{equation*}
        \begin{tikzcd}[row sep=0.1em]
            \mathllap{\Id_S \circ \Id_S \;\colon\:} s \arrow[r, maps to, "{\Id_S}"]
                & s \arrow[r, maps to, "{\Id_S}"]
                & s \arrow[d, equal] \\
            \mathllap{\phantom{\Id_S \circ{}}\Id_S \;\colon\:} s \arrow[rr, maps to, "{\Id_S}"']
                & {}
                & s
        \end{tikzcd}
    \end{equation*}
    The two functions agree on their assigned values, and therefore are equal. This covers both left and right compositions of the inverse, so $\Id_S$ is invertible with itself as its inverse.

    \medskip\noindent
    $\Id_S$ is injective if $\Id_S(s_1) = \Id_S(s_2)$ implies that $s_1 = s_2$. This is trivial:
    \begin{equation*}
        s_1 = \Id_S(s_1) = \Id_S(s_2) = s_2
    \end{equation*}
\end{example}
\begin{example}[Addition]
    Addition is normally called a binary operator, which is just a function on pairs:
    \begin{align*}
        +\, \colon \R \times \R &\longrightarrow\:\:\: \R \\
                        (a,b)\; &\longmapsto     a + b
    \end{align*}
    $+$ is not an injection since $1 + 2 = 0 + 3$ for example. ie. $+$ applied to the pairs $(1,2)$ and $(0,3)$ give the same result.

    \medskip\noindent
    $+$ is a surjection since any $r \in \R$ can be obtained as $0 + r$.
    Once we have proven Theorem \ref{bijective_invertible}, we'll see that this means that addition is not invertible.
\end{example}
\begin{example}[Addition by a constant]
    Addition itself is not an injection, but how about addition by a constant? Consider $(c+\phantom{})$ as a function $\R \to \R$ defined by
    \begin{equation*}
        (c+\phantom{}) \colon r \longmapsto c + r
    \end{equation*}
    Now, $(c + \phantom{})$ is injective if whenever $c + r_1 = c + r_2$, we have $r_1 = r_2$. But that's obvious because we can just subtract $c$ from both sides.

    This argument tells us exactly what we should expect the inverse of $(c + \phantom{})$ to be: $(-c + \phantom{})$. Indeed,
    \begin{align*}
        \big((c + \phantom{}) \circ (-c + \phantom{})\big)(r) &= (c + \phantom{})((-c + \phantom{})(r)) \\
        &= (c + \phantom{})(-c + r) \\
        &= c + -c + r \\
        &= r \\
        &= \Id_{\R}(r)
    \end{align*}
    This shows that $(c + \phantom{}) \circ (-c + \phantom{}) = \Id_{\R}$, we also need to show that $(-c + \phantom{}) \circ (c + \phantom{}) = \Id_{\R}$ which follows naturally by a similar chained equality. Thus, $(-c + \phantom{})$ is the inverse map to $(c + \phantom{})$.
    This is the way in which subtraction is the inverse of addition despite the fact that addition itself is not invertible.
\end{example}

\begin{theorem}[bijective if and only if invertible]\label{bijective_invertible}
    Let $f \colon A \to B$ be a map. $f$ is bijective if and only if it is invertible.
\end{theorem}
\begin{proof}
    Suppose $f$ is invertible, then it has inverse $g \colon B \to A$.

    \medskip\noindent
    Injective:

    \noindent
    Let $a_1, a_2 \in A$ be such that $f(a_1) = f(a_2)$, then $g(f(a_1)) = g(f(a_2))$, but $g$ is the inverse of $f$, so
    \begin{equation*}
        a_1 = (g\circ f)(a_1) = g(f(a_1)) = g(f(a_2)) = (g\circ f)(a_2) = a_2.
    \end{equation*}
    Surjective:

    \noindent
    Let $b \in B$, then $g(b) \in A$.
    \begin{equation*}
        f(g(b)) = (f \circ g)(b) = b
    \end{equation*}
    Thus, we have found an element of $A$ which maps onto $b$. This argument did not depend on $b \in B$, so any $b \in B$ has such an element of $A$.

    \medskip\noindent
    Conversely, suppose $f$ is bijective. We wish to construct a function $g \colon B \to A$ which is inverse to $f$.
    Let $b \in B$. Since $f$ is surjective, there is an $a \in A$ such that $f(a) = b$. Since $f$ is injective, only 1 such $a$ exists. This is true for any $b\in B$, so we can define:
    \begin{equation*}
        g \colon b \longmapsto \text{the unique element $a\in A$ such that $f(a) = b$}.
    \end{equation*}
    $g \circ f = \Id_A$:

    \noindent
    $(g \circ f)(a) = g(f(a))$ which is by definition, the unique $a_1 \in A$ such that $f(a_1) = f(a)$. That unique $a_1$ is $a$, so $g(f(a)) = a$.

    \medskip\noindent
    $f \circ g = \Id_B$:

    \noindent
    $(f \circ g)(b) = f(g(b))$. $g(b)$ is an element $a\in A$ with $f(a) = b$, so $f(g(b)) = b$.

    \medskip\noindent
    Together, we see that $g$ is indeed the inverse of $f$, so we are done.
\end{proof}
\begin{remark}
    When we look at the second implication's two pieces, proving $g \circ f = \Id_A$ used the uniqueness of $a_1$ (the injectivity property), while proving $f \circ g = \Id_B$ simply used that the value mapped onto $b$ (the surjectivity property).
    This gives some good indication that injectivity and surjectivity interact quite nicely on specific sides of compositions.
\end{remark}
